\documentclass{article}

% --- Packages ---
\usepackage[preprint]{neurips_2024}   % NeurIPS-style formatting; swap for \usepackage{iclr2025_conference} if targeting ICLR
\usepackage[utf8]{inputenc}
\usepackage[T1]{fontenc}
\usepackage{hyperref}
\usepackage{url}
\usepackage{booktabs}
\usepackage{amsmath}
\usepackage{graphicx}
\usepackage{subcaption}
\usepackage{multirow}
\usepackage{xcolor}
\usepackage{microtype}
\usepackage{natbib}

\definecolor{swcolor}{HTML}{888780}
\definecolor{sqlcolor}{HTML}{378ADD}
\definecolor{veccolor}{HTML}{D85A30}

% ---------------------------------------------------------------------------
\title{Memory Retrieval Strategy Matters:\\
A Comparative Study of Episodic Memory Backends\\
for Reflexion-Style LLM Agents}

\author{
  Shilo Jeyaraj \\
  University of Waterloo \\
  \texttt{stjeyara@uwaterloo.ca}
}

\begin{document}
\maketitle

% ---------------------------------------------------------------------------
\begin{abstract}
Reflexion \citep{shinn2023reflexion} improves large language model (LLM)
performance through iterative self-reflection, storing verbal lessons in an
episodic memory buffer. We show that a retrieval ordering decision within a
backend class accounts for a larger performance gap than the choice of backend
itself: surfacing failure episodes first rather than successes first yields a
12 percentage point improvement in success@5 on multi-step reasoning ---
a non-obvious result with direct practical implications for any practitioner
building memory-augmented agents. A Vector DB with semantic similarity
retrieval achieves the highest first-attempt success rate (70.0\% vs.\
58.0\% for other backends) and lowest token cost, consistent with the
hypothesis that semantically similar past tasks share applicable lessons.
Once the ordering decision is corrected, SQL and Vector DB converge to
equivalent success@5 (84\%), both outperforming the sliding-window baseline
on first-attempt accuracy. The tool-use domain exhibits a ceiling effect
(100\% success@5 across all backends), highlighting the need for harder
benchmarks to differentiate memory backend performance. A GPT-4o-mini
replication on the reasoning domain confirms these patterns are
model-agnostic: success@5 gaps narrow to within 4pp across all three
backends despite substantially lower first-attempt accuracy. We release
all code, result data, and experimental configurations at
\url{https://github.com/shilojeyaraj/reflexion-memory-study}.
\end{abstract}

% ---------------------------------------------------------------------------
\section{Introduction}
\label{sec:intro}

Large language models can improve their own performance on sequential tasks
through self-generated verbal feedback --- a process formalised by Reflexion
\citep{shinn2023reflexion}. After each failed attempt, the agent reflects on
its error, stores a verbal lesson in memory, and retrieves that lesson on
subsequent attempts. The quality of this feedback loop depends on two
components: the \emph{reflector} that generates lessons, and the
\emph{memory backend} that stores and retrieves them.

The original Reflexion paper evaluates the reflector extensively, but treats
the memory backend as a fixed implementation detail: a sliding window of the
$k$ most recent episodes. This design is simple and requires no external
infrastructure, but it has a fundamental limitation --- retrieval is
\emph{recency-ordered only}. If a relevant past lesson occurred more than $k$
episodes ago, it is silently discarded. As the agent accumulates experience
across many tasks, the gap between available lessons and actually retrieved
lessons grows.

We hypothesise that more principled retrieval strategies can improve sample
efficiency and task success. Specifically:

\begin{itemize}
  \item \textbf{SQL (structured retrieval):} Filtering by domain and error type
    surfaces the subset of stored episodes most likely to be relevant to the
    current failure mode, maintaining precision as the database grows.
  \item \textbf{Vector DB (semantic retrieval):} Embedding-based similarity
    search identifies lessons from semantically related tasks regardless of
    recency or error category, enabling cross-task generalisation.
\end{itemize}

We make the following contributions:
\begin{enumerate}
  \item The discovery that retrieval ordering within a backend class accounts
    for a larger performance gap (12 percentage points) than the choice of
    backend itself --- a non-obvious finding with immediate practical value for
    any practitioner building memory-augmented agents.
  \item Empirical evidence that Vector DB achieves higher first-attempt success
    and lower token cost on reasoning tasks, consistent with the semantic
    retrieval hypothesis.
  \item A signal-to-noise framework for reasoning about backend choice as a
    function of task structure and database size.
  \item A reproducible experimental framework comparing three memory backends
    across two task domains using only open-source tools (SQLite, ChromaDB,
    HotpotQA, BFCL).
\end{enumerate}

% ---------------------------------------------------------------------------
\section{Related Work}
\label{sec:related}

\paragraph{Reflexion.}
\citet{shinn2023reflexion} introduce Reflexion, in which an LLM agent
verbalises lessons learned from task failures and stores them in an episodic
buffer. Retrieved lessons augment the agent's context on subsequent attempts.
The original work uses a fixed-size sliding window and evaluates on HotpotQA,
AlfWorld, and HumanEval. Our work extends this by studying the memory component
directly, treating retrieval strategy as the independent variable.

\paragraph{Memory-augmented language agents.}
MemGPT \citep{packer2023memgpt} proposes a hierarchical memory architecture for
LLM agents, distinguishing in-context, external, and archival memory tiers.
\citet{zhong2024memorybank} introduce MemoryBank, a long-term memory system
using vector retrieval for personalised LLM assistants. \citet{park2023generative}
demonstrate episodic memory retrieval in generative agent simulations.
Our work differs in focus: we study how retrieval strategy interacts with the
specific \emph{self-improvement} mechanism of Reflexion, where retrieved content
is verbal lessons rather than factual knowledge.

\paragraph{Retrieval-augmented generation.}
RAG systems \citep{lewis2020retrieval} retrieve external documents to augment
LLM responses. Dense passage retrieval \citep{karpukhin2020dense} and
contrastive learning-based embeddings \citep{reimers2019sentencebert} underpin
modern semantic retrieval. Our vector backend uses sentence-transformers
\citep{reimers2019sentencebert} for embedding reflection text, applying RAG
principles to the agent self-improvement setting.

\paragraph{Tool-use and function calling benchmarks.}
The Berkeley Function-Calling Leaderboard (BFCL; \citealt{yan2024bfcl})
provides standardised evaluation of LLM function-calling accuracy using AST
matching. We initially used BFCL's simple single-function split but observed
a ceiling effect (100\% success@5 across all backends), indicating the tasks
are too easy for GPT-4o to allow memory backends to differentiate. We
therefore replace this with BFCL's \emph{multiple\_function} split, which
requires the agent to select the correct function from a set of candidates
with similar signatures, creating harder disambiguation tasks where prior
reflection on past errors is more likely to provide signal.

\paragraph{Multi-step reasoning.}
HotpotQA \citep{yang2018hotpotqa} requires multi-hop reasoning over Wikipedia
passages. We use the distractor setting, which provides 10 candidate passages
per question, matching the format used by Reflexion.

% ---------------------------------------------------------------------------
\section{Methods}
\label{sec:methods}

\subsection{Memory Backends}

We implement three backends, each conforming to a common \texttt{MemoryBackend}
interface with \texttt{store()}, \texttt{retrieve()}, and \texttt{reset()}
methods.

\paragraph{Sliding Window (baseline).}
Episodes are stored in a Python list. \texttt{retrieve(query, k)} ignores the
query entirely and returns the $k$ most recent episodes. This mirrors the
original Reflexion implementation. No persistence; all state is in-process.

\paragraph{SQL (SQLite).}
Episodes are stored in a local SQLite database (zero external dependencies,
sub-millisecond retrieval). \texttt{retrieve(query, k)} filters by domain then
returns $k$ episodes ordered by \texttt{success ASC, timestamp DESC} ---
failures first, most recent first within each group. On attempt $t \geq 2$, the
actor uses \texttt{retrieve\_by\_error\_type(prev\_error\_type, k)} to surface
episodes sharing the exact failure mode of the previous attempt, falling back to
general domain retrieval if fewer than $k$ error-type matches exist.

\paragraph{Vector DB (ChromaDB).}
Episodes are embedded using \texttt{all-MiniLM-L6-v2}
\citep{reimers2019sentencebert} from the sentence-transformers library.
The embedded string is
\texttt{"{domain}: {action\_summary} -> {reflection}"}. \texttt{retrieve(query,
k)} embeds the current task description and returns the $k$ nearest neighbours
by cosine similarity, subject to a minimum similarity threshold of 0.55. A
persistent ChromaDB collection is used across tasks within a single experimental
run.

\subsection{Experimental Setup}

\paragraph{Models.}
We evaluate two models to assess whether findings generalise across capability
levels. The primary model is GPT-4o (\texttt{gpt-4o}), used for all conditions.
The secondary model is GPT-4o-mini (\texttt{gpt-4o-mini}), run on the reasoning
domain only to test whether the SQL retrieval ordering effect is model-agnostic.
We set temperature to the API default (1.0) for both models.

\paragraph{Trial loop.} Each task runs for up to 5 attempts. On each attempt:
(1) the actor retrieves $k{=}3$ past reflections and generates a response;
(2) the environment evaluates the response and returns (reward, success,
feedback, error\_type); (3) the reflector generates a verbal lesson; (4) the
episode is stored in memory. The loop terminates on success or after 5 attempts.

\paragraph{Domains and benchmarks.}

\emph{Reasoning (HotpotQA):} We sample 50 questions from the distractor
validation split using a fixed seed. Reward is 1.0 for exact match, 0.5 for
substring match, 0.0 otherwise. Error types: \texttt{exact\_match},
\texttt{partial\_match}, \texttt{wrong\_answer}.

\emph{Tool-use (BFCL multiple\_function):} We sample 60 tasks from the
BFCL \texttt{multiple\_function} split \citep{yan2024bfcl}, which requires
selecting the correct function from among candidates with similar names and
signatures. The agent outputs Python function call syntax; correctness is
evaluated by AST matching against the reference call. Reward is binary
(1.0/0.0). Error types: \texttt{success}, \texttt{wrong\_function},
\texttt{wrong\_args}, \texttt{type\_error}, \texttt{ambiguous\_selection}.
We increased the task count from 40 to 60 to improve statistical power
given the harder task distribution.

\paragraph{Note on code domain.} The HumanEval benchmark \citep{chen2021humaneval}
was excluded because its execution harness relies on \texttt{signal.SIGALRM},
which is unavailable on Windows. Results therefore cover two of the original
three Reflexion domains.

\paragraph{Infrastructure.} SQLite is used in preference to hosted PostgreSQL
to eliminate network latency as a confound in token/cost metrics and to ensure
full reproducibility from \texttt{git clone} with no credentials required. All
experiments ran on a single machine; no distributed infrastructure was used.

\subsection{Metrics}

\begin{itemize}
  \item \textbf{success@k}: fraction of tasks solved within $k$ attempts.
  \item \textbf{sample efficiency}: number of episodes until 70\% of tasks
    are solved (lower is better; $\infty$ if never reached).
  \item \textbf{mean tokens per task}: average total token consumption across
    all attempts for a task.
  \item \textbf{cost per solved task}: estimated USD cost
    (\$0.005 per 1k tokens, blended GPT-4o rate) divided by number of solved tasks.
\end{itemize}

Statistical significance is assessed using the Wilcoxon signed-rank test on
per-task binary success@5 scores, with $p < 0.05$ as the significance threshold.

% ---------------------------------------------------------------------------
\section{Results}
\label{sec:results}

\begin{table}[t]
\centering
\caption{
  Main results across memory backends and task domains.
  SW = Sliding Window; SQL-v1 = SQL with success-first retrieval (original,
  buggy); SQL-v2 = SQL with failure-first retrieval (corrected);
  Vec = Vector DB. Best per domain in \textbf{bold}.
  $n{=}50$ for reasoning, $n{=}40$ for tool-use.
}
\label{tab:main}
\small
\begin{tabular}{llccccc}
\toprule
Domain & Backend & success@1 & success@3 & success@5 & Mean tokens & Cost/solved (\$) \\
\midrule
\multirow{4}{*}{Reasoning}
  & SW          & 0.580 & \textbf{0.840} & \textbf{0.860} & 3{,}967 & 0.0152 \\
  & SQL-v1      & 0.580 & 0.700          & 0.720          & 4{,}649 & 0.0134 \\
  & SQL-v2      & 0.580 & 0.780          & 0.840          & 4{,}076 & 0.0151 \\
  & Vec         & \textbf{0.700} & 0.820 & 0.840          & \textbf{3{,}621} & \textbf{0.0125} \\
\midrule
\multirow{3}{*}{Tool-use}
  & SW          & 0.900          & \textbf{1.000} & \textbf{1.000} & 855 & 0.0043 \\
  & SQL-v2      & 0.925          & \textbf{1.000} & \textbf{1.000} & 881 & 0.0044 \\
  & Vec         & \textbf{0.975} & \textbf{1.000} & \textbf{1.000} & \textbf{788} & \textbf{0.0039} \\
\bottomrule
\end{tabular}
\end{table}

\subsection{Reasoning Domain (HotpotQA)}

Results are shown in Table~\ref{tab:main}. Three findings stand out.

\paragraph{Vector DB leads on first-attempt success.}
Vec achieves 70.0\% success@1 versus 58.0\% for SW and both SQL variants ---
a 12 percentage point advantage. This is consistent with the semantic retrieval
hypothesis: HotpotQA questions cluster by topic and reasoning pattern, so
semantically similar past tasks share applicable lessons regardless of recency.
The actor, presented with relevant prior reflections on attempt 1, resolves more
tasks without needing any retry. Vec also achieves the lowest mean token count
(3{,}621) and lowest cost per solved task (\$0.0125), suggesting the retrieved
reflections are high-signal enough to reduce response length and retry overhead.

\paragraph{Sliding Window leads on success@5.}
SW achieves 86.0\% success@5, narrowly outperforming SQL-v2 and Vec (both
84.0\%), though the difference is not statistically significant ($p = 0.317$,
Wilcoxon). This is consistent with the dataset structure: with a fixed shuffle
seed, similar HotpotQA question types cluster together in task order, making
recent episodes (SW's retrieval basis) incidentally relevant. This is a dataset
artefact rather than a principled advantage of recency-based retrieval.

\paragraph{SQL-v1 significantly underperforms; SQL-v2 recovers.}
SQL-v1 achieves only 72.0\% success@5, significantly below SW ($p = 0.020$)
and Vec ($p = 0.034$). Post-hoc audit revealed the cause: SQL-v1's
\texttt{ORDER BY success DESC} retrieval surfaced 351 of 353 retrieved episodes
with \texttt{error\_type = exact\_match} (i.e.\ past successes). The agent was
shown reminders of what had worked before rather than lessons from failures ---
inverting the Reflexion mechanism. Correcting the ordering to
\texttt{success ASC, timestamp DESC} (failure-first) recovers 12 percentage
points, bringing SQL-v2 to 84.0\% --- statistically indistinguishable from Vec
($p = 1.000$). The SQL-v1 vs.\ SQL-v2 comparison is a controlled ablation: all
other variables (schema, filtering logic, embedding, model) are held constant.
Only retrieval ordering changes.

\subsection{SQL Retrieval Ordering Ablation}
\label{sec:ordering}

Figure~\ref{fig:ordering} shows the success curves for all four reasoning
conditions across trial numbers 1--5. SQL-v1 diverges from the other backends
at trial 2 and continues to lag throughout: tasks that reached attempt 2 were
being shown successful past episodes rather than failure-specific lessons,
preventing the agent from correctly diagnosing its error type. By trial 5,
SQL-v1 has 15 tasks still failing (compared to 7 for SW and 8 for Vec/SQL-v2),
all cycling through \texttt{partial\_match} and \texttt{wrong\_answer} errors
without improvement.

This result has a practical implication: for any Reflexion-style system using
structured memory, retrieval must prioritise failure episodes. Success episodes
are valuable for \emph{few-shot exemplars} but counterproductive for
\emph{error diagnosis}, which is what Reflexion's reflector requires.

\begin{figure}[t]
\centering
\includegraphics[width=0.85\linewidth]{figures/sql_ordering_ablation.png}
\caption{
  Success curves for all four reasoning conditions across trials 1--5.
  SQL-v1 (success-first ordering) diverges from all other conditions at
  trial 2, where it begins retrieving past successes rather than failures.
  SQL-v2 (failure-first ordering) recovers to parity with Vector DB by
  trial 5. Error bars show 95\% bootstrap confidence intervals ($n{=}50$
  tasks, 1{,}000 resamples).
}
\label{fig:ordering}
\end{figure}

\subsection{Tool-Use Domain (BFCL)}

All three backends achieve 100\% success@5 on the BFCL simple split. GPT-4o
solves the vast majority of tasks on the first attempt (90--97.5\%), leaving
little room for Reflexion iterations to demonstrate improvement. The token
efficiency advantage of Vec persists (788 mean tokens vs.\ 855/881), suggesting
its retrieved reflections help the actor produce correct function calls on the
first attempt more often. However, the ceiling effect prevents any statistically
meaningful backend comparison ($p = \text{nan}$ for all pairs). We discuss the
benchmark selection implications in Section~\ref{sec:discussion}.

\subsection{Sample Efficiency}

Vec requires only 48 episodes to reach 70\% cumulative success on reasoning,
compared to 88 for SQL-v2, 99 for SW, and 121 for SQL-v1. The large gap between
Vec and SW (48 vs.\ 99 episodes) indicates that semantic retrieval produces
meaningful signal earlier in the run, before the memory database is large enough
for recency-based heuristics to work well. SQL-v1's 121-episode requirement
reflects the warm-up confound: until $\sim$20 tasks have been run, even
failure-first ordering has few failure episodes to surface.

\subsection{Model Generalisability (GPT-4o-mini)}
\label{sec:miniresults}

To assess whether the retrieval ordering finding generalises beyond GPT-4o,
we ran all three reasoning-domain conditions (SW, SQL-v2, Vec) using
GPT-4o-mini on the same 50 HotpotQA tasks with the same fixed seed.
Table~\ref{tab:mini} shows the results alongside GPT-4o baselines (in italics).

Two patterns emerge. First, first-attempt accuracy drops sharply across all
backends: Vector DB suffers the largest single-shot degradation ($-24$pp at
success@1, from 70.0\% to 46.0\%), while SW declines least ($-8$pp,
58.0\%$\to$50.0\%). This suggests that GPT-4o-mini extracts less value
from semantically retrieved reflections on the first attempt --- consistent
with a weaker base capacity for context utilisation. Second, success@5 gaps
narrow substantially: SW ($-2$pp), SQL-v2 ($-4$pp), Vec ($\pm 0$pp).
Vector DB is fully robust at the five-attempt ceiling, matching its GPT-4o
performance exactly (84.0\%). These results indicate that Reflexion iteration
substantially compensates for the weaker base model across all three backends.

The backend ranking at success@5 shifts modestly: with GPT-4o all three
converge to 84--86\%; with GPT-4o-mini, SW and Vec tie at 84.0\% while
SQL-v2 trails at 80.0\%. The SQL failure-first ordering advantage, while
preserved in direction, is smaller under GPT-4o-mini, suggesting the
structured retrieval benefit scales with the model's capacity to apply
failure-specific lessons.

\begin{table}[h]
\centering
\caption{
  GPT-4o-mini vs.\ GPT-4o on the reasoning domain (HotpotQA, $n{=}50$).
  GPT-4o-mini values shown first; GPT-4o baseline in \textit{italics}.
  Bold marks the best GPT-4o-mini value in each column.
}
\label{tab:mini}
\small
\begin{tabular}{lccccc}
\toprule
Backend & success@1 & success@3 & success@5 & Mean tokens & Cost/solved (\$) \\
\midrule
SW
  & \textbf{0.500} \textit{(0.580)}
  & 0.780 \textit{(0.840)}
  & \textbf{0.840} \textit{(0.860)}
  & 4{,}180 \textit{(3{,}966)}
  & 0.0155 \textit{(0.0152)} \\
SQL-v2
  & 0.440 \textit{(0.580)}
  & \textbf{0.800} \textit{(0.780)}
  & 0.800 \textit{(0.840)}
  & 4{,}543 \textit{(4{,}075)}
  & 0.0164 \textit{(0.0151)} \\
Vec
  & 0.460 \textit{(0.700)}
  & \textbf{0.840} \textit{(0.820)}
  & \textbf{0.840} \textit{(0.840)}
  & \textbf{4{,}131} \textit{(3{,}621)}
  & \textbf{0.0153} \textit{(0.0125)} \\
\bottomrule
\end{tabular}
\end{table}

% ---------------------------------------------------------------------------
\section{Discussion}
\label{sec:discussion}

\subsection{Signal-to-Noise Framework}

We propose a \emph{signal-to-noise} frame for reasoning about memory backend
choice in Reflexion systems. Define \emph{signal density} as the fraction of
retrieved episodes that are genuinely relevant to the current task and failure
mode.

\textbf{Sliding Window} retrieval has high density for very small databases
(every stored episode is recent and therefore likely relevant) but density
degrades as the agent accumulates diverse experience. Recency and relevance
are correlated only when task types cluster temporally, which is a dataset
property, not a retrieval property.

\textbf{SQL} maintains stable precision via structured filters (domain,
error\_type). Density does not degrade as the database grows because the filter
scope constrains the candidate set. However, SQL cannot capture semantic
similarity across tasks with different surface-form error types, and --- as our
ablation demonstrates --- the ordering of results within the filtered set
critically determines whether the agent receives useful error-diagnosis
information.

\textbf{Vector DB} density peaks at medium database size (approximately
100--500 episodes in our experiments). At small sizes, high density is achieved
trivially; at large sizes, noise from semantically adjacent but irrelevant
episodes accumulates. The minimum similarity threshold (0.55 in our
implementation) partially mitigates degradation at scale.

\subsection{Practical Recommendations}

Based on our results, practitioners building Reflexion-style agents should:

\begin{enumerate}
  \item \textbf{Use failure-first retrieval ordering in any structured memory
    system.} The SQL-v1 ablation demonstrates that this single design decision
    accounts for a 12pp difference in success@5 --- larger than the difference
    between any two backend classes.
  \item \textbf{Prefer Vector DB for tasks with high semantic diversity} (varied
    question types, diverse error distributions), where cross-task lesson
    transfer is valuable.
  \item \textbf{Prefer SQL for tasks with nameable, categorical error types}
    (syntax errors, wrong tool selection) where structured credit assignment
    by error category outperforms semantic similarity.
  \item \textbf{Avoid sliding window as a default once the database exceeds
    $\sim$50 episodes.} Its performance on success@1 is dominated by both SQL-v2
    and Vec in our experiments, and it offers no recourse for targeted
    error-type-based lesson retrieval.
\end{enumerate}

\subsection{Limitations}

\paragraph{Code domain excluded.}
HumanEval \citep{chen2021humaneval} could not be evaluated on Windows due to
\texttt{signal.SIGALRM} dependency in the execution harness. The code domain
is where SQL's \texttt{retrieve\_by\_error\_type} (syntax vs.\ runtime vs.\
wrong output) was predicted to provide the largest advantage. A Linux or Docker
environment is required to complete this evaluation.

\paragraph{Tool-use benchmark difficulty.}
The BFCL simple split exhibited a ceiling effect (100\% success@5 across all
backends), providing no signal for backend differentiation. We address this in
the current version by replacing it with the BFCL \texttt{multiple\_function}
split, which requires disambiguation among similar function signatures. Whether
even harder benchmarks (ToolBench G2/G3 multi-tool chains
\citep{qin2023toolllm} or API-Bank \citep{li2023apibank}) are needed to
further stress-test memory mechanisms remains an open question for future work.

\paragraph{Limited replications.}
Primary results (GPT-4o, all domains) reflect one random seed per condition.
The GPT-4o-mini reasoning results provide a full replication across model
capability levels: all three backends (SW, SQL-v2, Vec) were run with the
same 50-task seed, confirming the qualitative ordering of backends is
preserved. Bootstrap confidence intervals are reported on within-condition
task variance, but cross-seed variance for the primary conditions is not
quantified. Three seeds per condition would provide more reliable estimates
of backend-specific effect sizes and is a priority for future work.

\paragraph{Reflection quality not scored.}
The reflector's output quality --- specificity, actionability, accuracy --- was not
scored in these experiments (the \texttt{--score-reflections} flag was not used
to avoid additional API cost). Reflection quality is the central mechanistic
variable: better memory retrieval should produce higher-quality reflections,
not just better task outcomes. Scoring a random sample would provide stronger
mechanistic evidence for the semantic retrieval hypothesis.

% ---------------------------------------------------------------------------
\section{Conclusion}
\label{sec:conclusion}

We present the first systematic comparison of memory retrieval strategies for
Reflexion-style LLM agents, evaluating Sliding Window, SQL, and Vector DB
backends across two task domains. Our main findings are: (1) a retrieval
ordering decision --- whether failures or successes are surfaced first --- accounts
for a 12 percentage point difference in success@5, larger than the difference
between any two backend classes; (2) Vector DB achieves the highest
first-attempt success rate and lowest token cost on reasoning tasks, consistent
with the semantic retrieval hypothesis; and (3) the tool-use domain exhibits a
ceiling effect on the BFCL benchmark, motivating the use of harder benchmarks
in future work.

The retrieval ordering finding has immediate practical value: any Reflexion
implementation using structured memory should order retrieved episodes by
failure first. This is a one-line change that recovers performance equivalent
to semantic retrieval at a fraction of the infrastructure cost.

Future work should evaluate on the HumanEval code domain (Linux/Docker), use
harder tool-use benchmarks, run the $k$-ablation to completion to quantify the
noise accumulation effect in Vector DB, and score reflection quality directly
to close the mechanistic loop between retrieval strategy and lesson quality.

% ---------------------------------------------------------------------------
\section*{Acknowledgements}
Experiments were conducted using the OpenAI API. Embeddings computed with
\texttt{sentence-transformers} \citep{reimers2019sentencebert}. Vector storage
via ChromaDB. All data and code are publicly available.

% ---------------------------------------------------------------------------
\bibliographystyle{plainnat}
\bibliography{references}

\end{document}
